\documentclass[11pt,a4paper]{article}
\usepackage[margin=24mm]{geometry}
\usepackage{amsmath,amssymb,booktabs,longtable,array,hyperref}
\hypersetup{hidelinks}
\newcommand{\Lc}{\textbf{(c)}\ }
\newcommand{\Ld}{\textbf{(d)}\ }
\newcommand{\scr}{\texttt{audit\_smx.py}}
\newcolumntype{L}[1]{>{\raggedright\arraybackslash}p{#1}}
\emergencystretch=3em
\title{Independent audit of the S-matrix candidates package}
\author{Quantum Boussinesq solitons project}
\date{4 November 2026}
\begin{document}
\maketitle
\begin{abstract}
\Lc The classical dictionary, conditional Gamma algebra, selected breather poles,
and a directional phase derivative survive independent checks. Several printed
claims do not: Darboux signs, a resonance line, the breather positivity threshold,
the continued antisymmetric continuum integral, and the small-amplitude frame
formulas. The four original scripts nevertheless all print
\texttt{ALL CHECKS PASSED}. The quantum-dilogarithm exclusion is too broad;
a narrower finite root-ratio exclusion has a repairable proof.
\Ld Neither quantum ring holonomy nor a unique physical one-loop phase follows
from the available consistency identities. The frame dictionary's antisymmetric
two-body freedom remains an independent input to the proposed interface.
\end{abstract}

\section{Scope, labels, and independent implementation}
Labels mean (a) literature proof, (b) literature heuristic, (c) project
derivation or explicitly numerical evidence, and (d) guess or unresolved
physical interpretation. Verdicts apply to the specified statement, not to
an entire research programme. All mathematical formulas advanced as (c) here
are checked in \texttt{scripts/}\scr; section names below identify its functions.
Numerical checks support the displayed algebra and counterexamples; finite
sampling is not represented as a global proof.

\Lc Inputs read were \texttt{AGENTS.md}, the status in \texttt{README.md},
\texttt{audit\_phase2i.tex} as the format and dependency audit,
\texttt{smx\_candidates.tex}, text extracted from its 18-page PDF, all four
\texttt{smx\_*} scripts, \texttt{frame\_dictionary.tex}, and the relevant
breather and integral definitions in \texttt{phase1/review.tex},
\texttt{oneloop\_phase.tex}, and the earlier audit checker.
No external paper or abstract was read for this audit and no new external
reference is cited. In particular, attribution to Bogdan--Charkina below means
comparison with the action formula recorded in the local project texts,
not an independent verification of their paper or an exact quantum spectrum.
The review itself identifies its good-equation no-Hirota-breather argument
as a project result (c), not an established literature theorem (a).

\Lc The independent checker imports no project code. It reconstructs the
interaction from KP chords, checks it separately against the Hirota polynomial,
differentiates the classical action, and integrates the elementary phonon
phase on both open half-lines. SymPy and 60-digit mpmath are used. Exact-identity
checks use absolute tolerance $10^{-48}$; finite-step asymptotic checks use
explicit looser tolerances appropriate to their step sizes. Original scripts
run only in separate subprocesses for the requested execution audit, with
bytecode writing disabled; none supplies an independent-computation target.
The checker hashes the candidate TeX, PDF, and scripts before and after running.
Only the requested new files are delivered. No git operation, existing-file
edit, physical convention choice, or unrelated full-suite run was performed.

\clearpage
\section{Verdict table}
\small
\begin{longtable}{@{}L{5.0cm}L{2.1cm}L{8.0cm}@{}}
\toprule Claim & Verdict & Independent evidence / limit\\\midrule\endhead
Good/bad charge dictionary, KP ratio, log derivative & Confirmed (c) & \scr: \texttt{classical}; both equations and four directions.\\
Gamma shifts, Stirling terms, displayed CDD asymptotics; coupling reparametrisation & Confirmed (c) & \texttt{families}; conditional algebra on consistent branches, not physical selection.\\
Rational mixed derivative of classical phase & Confirmed (c) & Exact rational uniformisation and symbolic differentiation.\\
Universal mixed-argument quantum-dilogarithm exclusion & Refuted (c) & Arbitrary-argument counterexample and unproved cancellation premise; a generic single algebraic block obstruction survives.\\
Nine KP-ratio dilogs with constant coefficients & Qualified & Exclusion confirmed (c) by corrected residue argument; printed residue and proof are wrong. The 0.0368 fit is evidence only.\\
All mirror products excluded by unitarity alone & Qualified & Smooth nonconstant $G$ excluded; constant $G$ and exchange-unitary $F$ allowed. Both separated sides must be constant.\\
Bad Gamma poles, no $N=0$ pole & Confirmed (c) & Full-candidate residue and threshold tests on the stated breather window. Exact physical spectrum remains (d).\\
Good $A<1/4$, singular two-phase zero-background breathers & Confirmed (c) & Independent Hirota evaluation and tau positivity.\\
Good sign boundary, regularity threshold, numerical pole examples & Refuted (c) & Boundary is $a=\pi/3$; positivity requires $A>1$; corrected pole numbers below.\\
Good Gamma poles are spurious physical states & Qualified & Candidate poles exist (c); physical-sheet/state identification is an imposed assumption (d).\\
All three resonance lines and smoothness & Refuted (c) & $|D|=\pi/3$ is not a resonance; actual limits are finite but generally not differentiably smooth.\\
Darboux chamber form and seam endpoint difference & Confirmed (c) & Wedge algebra for $N=2,3,4$, all chambers; endpoint difference follows from coordinates.\\
Printed Darboux wall signs and total Bethe derivative & Refuted (c) & Opposite signs; omitted self-winding derivative.\\
Directional ring derivative, sum, composite charge identities & Confirmed (c) & \texttt{classical,poles}; conditional on the defined twist and proposed equations.\\
Quantum ring twist / exact field-theoretic quantisation & Qualified & Generating-function consistency is not a derivation of quantum holonomy (d).\\
Single overtaking residual, crossing-locked head-on channel & Qualified & Local analytic phase parametrisation only; additional zero-free, branch, and expansion assumptions needed.\\
One-body frame inputs uniquely fix residual & Refuted (c) & Frame dictionary's explicit antisymmetric two-body counterexample survives.\\
Raw continuum $I$ crossing & Confirmed (c) & Direct phonon quadrature; applies to the earlier definitions, not the malformed printed formula.\\
$J$ crossing sum $+\pi/2$ and cusp/Stirling cancellation & Refuted (c) & Actual antisymmetrised integral has crossing sum zero; checker inserts the offset.\\
Printed small-amplitude interface & Refuted (c) & Wrong tails, mode coefficient, head-on expansion, and $J\to0$; corrected formulas below.\\
Phase 2i independence; scripts terminate successfully & Confirmed (c) & Source inspection, import inspection, and subprocess execution.\\
\bottomrule
\end{longtable}
\normalsize

\section{Classical algebra and the family exclusions}
\paragraph{Dictionary and Gamma algebra (\texttt{classical,families}).}
\Lc With the stipulated charges and the independently reconstructed $A$,
\begin{align}
 P'_1&=\frac{2k_1}{9}\ell'_1,\\
 \frac29(\ell_1-\ell_2)\partial_1\log A
 &=-\frac13\Pi'_1M_2-\frac23M'_1\Pi_2,\\
 \partial_1\delta&=\frac{P'_1}{k_1}\log A-\Pi'_1M_2.
 \label{eq:derivative}
\end{align}
These hold in both equations and all tested directions. Charge-energy
differentiation also gives $dE/dP=v$. The KP root cubics and crossing of both
the dynamical and topological action pieces are confirmed independently.

\Lc The shifted Gamma quotient has offset
$\tau(s-\tfrac12)\pi\operatorname{sgn}\eta$ and next term
$-\tau(s^2-s+\tfrac16)/\eta$, with $\tau=+1$ for overtaking and $-1$ for
$+-$ head-on. The displayed rational/hyperbolic CDD examples have the stated
constant-step limits for fixed parameters. This is not a theorem about
\emph{every conceivable} function of $\eta$, particularly if its parameters
scale with $\hbar$. Global versus dynamical-only coupling reparametrisation
shifts the one-loop phase by $-c_1\delta$ versus $-c_1\delta_{\rm dyn}$.
\Ld These are reparametrisation identities, not a renormalisation prescription.
The family table's dynamical-only entry should retain this distinction.

\paragraph{The weight argument: a valid calculation with an invalid universal conclusion.}
\Lc Put $L=\log R$. Direct differentiation gives
\begin{equation}
 \partial_{12}\operatorname{Li}_2(R)
 =\frac{R}{1-R}L_1L_2-\log(1-R)L_{12}.
 \label{eq:dilog}
\end{equation}
Meanwhile $\partial_{12}[(\ell_1-\ell_2)\log A]$ has no undifferentiated
logarithm. In exponential rapidities it is rational, as is the derivative
of the topological term. For a generic single \emph{algebraic} nonseparable
argument such as $A$, the surviving logarithm obstructs a match. For
$R=f_1(\beta_1)/f_2(\beta_2)$ it does not: $L_{12}=0$.
This derivative diagnostic does not by itself exclude such separable arguments.

\Lc The text says ``any twice-differentiable argument'', and then extends
the single-block argument to arbitrary linear combinations without proving
independence of their logarithms. Neither extension follows. For a concrete
arbitrary-argument/product counterexample, locally for real $x,y$,
\begin{equation}
 \operatorname{Li}_2(-e^{xy})+\operatorname{Li}_2(-e^{-xy})
 =-\frac{\pi^2}{6}-\frac{x^2y^2}{2},\qquad
 \partial_{xy}(\text{left side})=-2xy.
\end{equation}
Both arguments have nonzero mixed logarithmic derivative, but their sum has
rational mixed derivative. This does not construct a Boussinesq S-matrix;
it refutes the proposed universal exclusion mechanism. Restricted rational
argument sets require their own functional-independence or singularity proof.

\paragraph{The finite KP-ratio basis has a repaired exclusion.}
\Lc Let $d=\alpha_1-\alpha_2$, $\alpha_2=\alpha$, and
$\varphi_r=\alpha+2\pi r/3$. The correct diagonal residues are
\begin{equation}
 \partial_{12}\operatorname{Li}_2(R_{rr})
 =\frac{1}{d\sin\varphi_r}+O(1),\qquad
 \partial_{12}\delta_{\rm dyn}^{\rm good}
 =-\frac{4\sin3\alpha}{3d}+O(1).
 \label{eq:residue}
\end{equation}
The bad classical residue is $-4\sinh3\beta/(3d)$.
The printed $-1/(2\sin\varphi_r)$ is wrong in sign and factor, and the
printed classical residues have the wrong sign. Moreover $1/\sin\varphi_r$
is unbounded at a zero of its denominator; the stated boundedness argument
is false.

\Lc Nevertheless a constant-coefficient sum of the three cosecants cannot
equal $-4\sin3\alpha/3$ on an open interval. By meromorphic continuation,
inspect $\alpha=-2\pi r/3$ modulo $\pi$: exactly one of the three cosecants
has a pole there, whereas the target is analytic. Its coefficient must vanish.
Doing this for each $r$ forces all three coefficients to vanish. Off-diagonal
root pairs and $\partial_{12}T$ have no generic diagonal $1/d$ pole and cannot
repair it. Thus the specific nine-ratio basis, with constant coefficients
and the allowed topological column, is excluded. This does not exclude
arbitrary ratios, variable prefactors, or general quantum-dilogarithm products.
\Lc Running the original fit reproduces 0.0368; its 40 float64 sample points
are evidence for that finite basis, not a proof. Our exclusion instead uses
the independently verified pole residues and the analytic argument above.

\paragraph{Mirror products: both separated sides must be constant.}
\Lc On a connected open rapidity domain, with nonzero differentiable factors,
\begin{equation}
 F(D)F(-D)G(S)^2=1
 \quad\Longrightarrow\quad F(D)F(-D)=C,\quad G(S)^2=C^{-1}.
\end{equation}
Thus $G$ is constant on each connected smooth branch. Independence of $S,D$
is valid locally even on the wedge of positive rapidities. The user's
suggested intermediate equation is exactly the right argument; it is
\emph{possible} when both sides are constant. For example $G=1$ and
$F(D)=e^{iD^3}$ satisfy exchange unitarity. Unitarity excludes a nontrivial
smooth mirror factor, not every mirror product. Without continuity one could
choose sign-changing $G(S)\in\{+1,-1\}$; regularity is an assumption.
A factor of a larger amplitude need not itself obey exchange unitarity
unless its complement does. The separate semiclassical nonseparability
obstruction remains valid for the proposed $F(D)G(S)e^{iT/\hbar}$ ansatz.

\paragraph{Exponential family: qualified, not generally entire.}
\Ld Calling an exponential pole-free requires a finite holomorphic exponent.
Here $\delta$ contains logarithms of $A$, and the correction is unrestricted;
neither entire dependence on $\eta$ nor global pole freedom has been proved.
\Lc Already $\exp(i[i\log z])=1/z$ locally illustrates why ``exponential''
alone does not rule out poles. \Ld Absence of classical breathers does not
establish absence of all quantum bound states; conversely a semiclassical
breather action alone does not establish an exact quantum pole requirement.
The physical family-V selection is therefore conditional.

\section{Breathers, poles, and resonance claims (\texttt{poles})}
\paragraph{Bad equation: confirmed candidate poles, qualified physical interpretation.}
\Lc For $\beta_{1,2}=a\mp ib$, $a>0$, $\pi/3<b<\pi/2$,
\begin{equation}
 A=\frac{1+3\coth^2a}{1-3\cot^2b}>1,\qquad
 i\eta=\frac{4\sinh3a\sin3b}{9\hbar}=-\frac{I_{\rm br}}{\hbar}.
\end{equation}
The relative momentum has positive imaginary part. Independent explicit
inversion of this action locates levels $N=1,2,3$ at $a=0.6,0.9$ and
$\hbar=0.1$. Multiplying the \emph{full candidate} by displacement from the
pole gives a finite nonzero residue, consistent across displacements
$10^{-15},10^{-19}$. Half-shifted Gamma poles move by a half-integer as
claimed. At fixed $a$ only levels within the action range of the specified
window are realised. No pole occurs at $N=0$: the full candidate divided
by its topological factor tends to one. This threshold limit can contain
$t\log t$ terms and need not be analytic. The original checker's comment
that $\Gamma(i\eta)/\Gamma(-i\eta)$ has a zero-level pole is also wrong:
its removable limit is $-1$.
\Ld Matching these candidate poles to the recorded Bohr--Sommerfeld action
is not an exact quantum-spectrum theorem. The inequality $A>1$ itself also
holds at $b=\pi/2$; excluding that endpoint for localisation is a separate
condition, not the printed algebraic equivalence.

\paragraph{Good equation: correct bound, incorrect sign partition and positivity test.}
\Lc Direct Hirota evaluation confirms
\begin{equation}
 A=\frac{1-3\cot^2a}{1+3\coth^2b}<\frac14.
\end{equation}
Its sign changes at $a=\pi/3$, not $\pi/6$: it is negative below
$\pi/3$ and positive above. Excluding $a=\pi/6$ has no stated basis.
Writing $X=e^R>0$, the real tau function is
\begin{equation}
 \tau=1+2X\cos I+AX^2.
\end{equation}
Its minimum over $I$ is $1-2X+AX^2$, which stays positive for all $X>0$
iff $A>1$. For example $A=1/2>1/4$, $X=2$, $I=\pi$ gives $\tau=-1$.
This is exactly the threshold in the Phase 1 review. Therefore the
zero-background two-phase Hirota class is singular, but the printed
positivity argument and attribution are wrong. \Ld A statement about every
possible good-equation breather, including other backgrounds, exceeds this
result and the review's explicitly limited scope.

\Lc The two momentum/spectral formulas in the pole proposition are correct:
\begin{equation}
 \operatorname{Im}(P_1-P_2)=\frac{\sinh2b}{3}
 (\cos4a\cosh2b-\cos2a),\qquad
 i\eta=\frac{4\sin3a\sinh3b}{9\hbar}.
\end{equation}
The stated inequality on $b$ gives examples with positive imaginary relative
momentum and $i\eta=-N$. At $N=1$, $\hbar=0.05$ the actual numbers are
\begin{center}
\begin{tabular}{rrrr}\toprule
$a$ & $b$ & $\operatorname{Im}(P_1-P_2)$ & $A$\\\midrule
0.4 & -0.040137301855 & 0.019447334601 & -0.008461762797\\
0.65 & -0.040269708896 & 0.030293216242 & -0.002261837165\\\bottomrule
\end{tabular}
\end{center}
Both $A$ values are negative, contrary to the printed second example.
The complete candidate has nonzero simple residues at these points on a
locally chosen logarithm branch. \Ld They are unwanted poles under the
package's spectrum criterion; calling this the physical quantum sheet uses
that criterion rather than deriving it from a quantum spectral problem.
For small $a<\pi/8$ not every negative $b$ satisfies the sheet inequality.

\paragraph{Resonance locus and regularity: refuted as printed.}
\Lc In the real square the genuine lines stated in the package are
$\alpha_1+\alpha_2=\pi/3$ (head-on) and
$\alpha_1+\alpha_2=2\pi/3$ (overtaking). The proposed
$|\alpha_1-\alpha_2|=\pi/3$ overtaking line is false: the overtaking
denominator vanishes at $|D|=2\pi/3$, outside this square.
At $\alpha_1=0.2$, $\alpha_2=0.2+\pi/3$ one finds
$A=0.355465479849$ and $\ell_1-\ell_2=-1.65067122982$, not infinity and zero.
The two genuine lines do give the finite limit $S\to e^{iT/\hbar}$.
However, approaching transversely by $t$ produces a nonzero $t\log t$
term; the logarithmically growing difference quotient is checked at
$t=10^{-10},10^{-20}$. ``Smoothly'' must be replaced by ``continuously,
with no pole in this limit'', not differentiably smooth.

\section{Ring geometry: established algebra and missing quantum input}
\paragraph{Darboux signs (\texttt{classical}).}
\Lc Starting from the stated chamber coordinate
$X'_j=X_j+\gamma_j\sum_{l\text{ left of }j}M_l$,
a particle moving from left to right gains the term $\gamma_jM_l$.
Consequently the consistent wall shifts are
\begin{equation}
 \Delta X'_j=\Delta x_{jl}+\gamma_jM_l,
 \qquad \Delta X'_l=\Delta x_{lj}-\gamma_lM_j.
 \label{eq:correctwall}
\end{equation}
They agree with $-\partial\delta/\partial P$ in
\eqref{eq:derivative}. The printed wall equation, introductory definition,
and frame split have the opposite signs. In contrast the residual-interface
\emph{script} uses the corrected signs, explaining why it passes.
The wedge identity $dP_j\wedge d(\gamma_jM_l)=d\Pi_j\wedge dM_l$
is correct since $d\gamma_j$ is proportional to $dP_j$. The change of ordering
and the Hirota shifts preserve the stated chamber two-form.

\paragraph{What the seam formula proves.}
\Lc With the package's explicit order of subtraction, the endpoint difference is
\begin{equation}
 X'_j(\Lambda^-)-X'_j(0^+)=\Lambda+\gamma_j\sum_{l\ne j}M_l.
\end{equation}
It is the negative of the instantaneous coordinate reset from $\Lambda^-$
to $0^+$. Calling both the seam jump obscures this sign. At fixed other charges,
the function $P_j\Lambda+\Pi_j\sum_{l\ne j}M_l$ has this derivative with
respect to $P_j$. It is obtained by integrating the endpoint difference
\emph{against $dP_j$}, not by the text's $P_j\,d(\Delta X'_j)$.
These two differential expressions are not equal.

\Ld To turn that generating-function observation into a quantum Bethe twist
requires a specified identification of the winding potential, its zero-mode
state and canonical one-form, and transition phases for all coordinates.
The package gives no wavefunction transport calculation establishing those
choices. In its own two-form the variable $c$ couples to $\Pi_{\rm tot}$,
not just $\Pi_j$. Extracting a self term $\Pi_jM_j$ from a seam change of $c$
requires keeping the simultaneous changes of the other coordinates.
Moreover the abstract says $c$ drops whereas the body says it increases.
We do not select a global boundary condition to fill this gap. The proposed
Aharonov--Bohm interpretation and exact quantum Bethe quantisation remain
heuristic (d); asymptotic line-soliton coordinates are not an exact finite-ring
field construction.

\paragraph{Directional derivative, sum rule, fusion: conditional confirmation.}
\Lc Define the proposed phase $\Phi(j\mid l)=\delta(j,l)+\Pi_jM_l$.
Then direct differentiation gives
\begin{equation}
 -\partial_{P_j}\Phi(j\mid l)
 =\Delta x_{jl}+\gamma_jM_l-\gamma_jM_l
 =-\frac{\log A_{jl}}{k_j}.
\end{equation}
Thus the final identity is right despite the intermediate sign errors.
It holds locally away from $P'_j=0$ and singular $A$. For the \emph{full}
classical Bethe left side there is, however, an omitted derivative:
\begin{equation}
 \partial_{P_j}\mathcal B_j
 =\Lambda-\sum_{l\ne j}\Delta x_{jl}
   +\partial_{P_j}(\Pi_jM_j).
\end{equation}
At parameter $0.23$, direction $+$, that last term is
$-17.3832401711$ (bad) and $-17.3822877701$ (good), not zero.

\Lc Pairwise exchange antisymmetry gives
$\Phi(j\mid l)+\Phi(l\mid j)=\Pi_jM_l+\Pi_lM_j$.
Summing the \emph{assumed} equations therefore gives
\begin{equation}
 P_{\rm tot}\Lambda+\Pi_{\rm tot}M_{\rm tot}
 =2\pi\hbar\sum_j n_j.
\end{equation}
For a quantum amplitude the unwrapped phases cancel modulo $2\pi$;
integer branch changes are absorbed in the $n_j$. This is an exact algebraic
consequence of the proposed Bethe equations, not an independent field-theory
derivation of those equations. The checker verifies the sum for $N=2,3,4$.

\Lc For the bad conjugate pair the sums
$M_{\rm br}=4\sinh a\cos b$ and
$\Pi_{\rm br}=-2\sinh2a\cos2b$ are real, and the fused third-particle
phase differentiates to $-\log(A_{31}A_{32})/k_3$ on conjugate logarithm
branches. This establishes charge additivity and a continued classical
shift. \Ld It does not prove a quantum bootstrap, normalised bound-state
residue, or bound-state ring Hilbert space.

\section{Residual phase and the frame dictionary (\texttt{interfaces})}
\Lc Locally on a real regular channel, two nonzero unitary amplitudes with a
common leading phase differ by $e^{i\Theta}$, with a real phase modulo
$2\pi$. Exchange imposes antisymmetry with the species labels exchanged.
If both amplitudes obey the stipulated crossing relation on compatible
analytic branches, then
\begin{equation}
 \Theta_{\rm ho}(\beta_1,\beta_2)
 =-\Theta_{\rm ov}(\beta_1,\beta_2+i\pi)\pmod{2\pi}.
\end{equation}
Thus a \emph{specified analytic} overtaking phase determines its crossed
channel. An arbitrary real antisymmetric function need not admit that
continuation or produce a real head-on phase. This is a functional freedom,
not one free real number.

\Ld Constraints (i)--(v) do not prove a globally zero-free ratio, a
single-valued holomorphic logarithm, or an integer-power loop expansion.
Zeros and monodromy are not eliminated merely by listing poles.
\Lc A nonzero entire multiplier with an antisymmetric phase proportional to
$\hbar^{-1/2}$ still has vanishing contribution to
$\hbar\arg S$ in the classical limit, while violating the asserted
$\Theta=\vartheta+O(\hbar)$ expansion. For example take the additional phase
$\hbar^{-1/2}\sinh(\beta_1-\beta_2)$ in both channels; it is real and
antisymmetric, and crossing reverses the hyperbolic sine. The claim needs an additional perturbative
expansion assumption. The quoted strip in $\operatorname{Im}(-i\eta)$
is not justified as a physical analyticity domain by these tests.

\paragraph{Relation to the frame dictionary.}
\Lc The frame dictionary exhibits
\begin{equation}
 A_\lambda(y_1,y_2)=\lambda\frac{y_1y_2(y_1-y_2)}{(y_1+y_2)^3},\qquad
 A_\lambda(Y,r^2Y)=\lambda\frac{r^2(1-r^2)}{(1+r^2)^3}.
\end{equation}
Our symbolic check independently verifies antisymmetry, edge vanishing,
nonzero mixed derivative, and the nonzero scale-invariant limit. Adding
$hA_\lambda$ to its action representative preserves one-body dispersion,
its classical limit, and symmetric reciprocity, but changes the phase by
$A_\lambda$. It therefore supplies precisely the same \emph{kind of}
antisymmetric two-body obstruction that can contribute to
$\vartheta_{\rm ov}$. It is not a one-body reparametrisation.

\Ld They are not already identified as the same function: the frame term is
one possible contribution to the total residual relative to Gamma, and its
finite-sea variables still need a physical continuum map. A global
crossing-real, pole-compatible extension of the displayed frame example
is not established by the frame dictionary. Nevertheless its existence
refutes the package's assertion that one-body $\Delta P$, a coupling shift,
and an intrinsic determinant uniquely determine the residual: a genuinely
two-body quantum frame input remains missing. The corrected interface would
need that contribution in addition to the one-body chain-rule terms. Neither
crossing nor the summed ring equation computes it.

\section{Direct one-loop quadrature: crossing and the cusp}
\paragraph{Definitions actually integrated (\texttt{quadrature}).}
\Lc The earlier note defines
\begin{equation}
 I_{12}=\frac1{4\pi}\int_{\mathbb R\setminus\{0\}}\theta_1\theta'_2\,dq,
 \qquad J_{12}=\frac{I_{12}-I_{21}}2,
\end{equation}
with its three-root elementary phase and weights $(1,1,-2)$.
We integrate that phase and its independently differentiated derivative,
splitting at zero and at each Lorentzian width. These are two choices
sharing a local gradient up to endpoint constants, not two contributions
to add. The new note's $\Delta\theta_{\rm cont}=I+J$ is not the earlier
one-loop prescription. Its displayed dilogarithm expression instead uses
$(1,1,-1)$ and a different kernel; its arctangent expression for $J$ is
also different. The original checker does not implement either printed
expression. Direct evaluation asserts discrepancies for both.

\Lc Under $\alpha\mapsto\alpha+\pi$, the elementary principal phases obey
$\theta(q;\alpha+\pi,+)=-\theta(q;\alpha,-)$, including their derivatives
on the open half-lines. Consequently
\begin{equation}
 I^{\rm same}_{12}(\alpha_1,\alpha_2+\pi)
 =-I^{\rm head}_{12}(\alpha_1,\alpha_2),\qquad
 J^{\rm same}_{12}(\alpha_1,\alpha_2+\pi)
 =-J^{\rm head}_{12}(\alpha_1,\alpha_2).
 \label{eq:truecross}
\end{equation}
The first requested identity is confirmed for the correct raw integral;
the second requested sum equals zero, not $\pi/2$. Numerical values are
\begin{center}\small
\begin{tabular}{rrrrr}\toprule
$(\alpha_1,\alpha_2)$ & $I^{\rm same}$ & $J^{\rm same}$ & $I^{\rm head}=J^{\rm head}$ & $J_{\rm cross}+J^{\rm head}$\\\midrule
$(0.2,0.55)$ & -0.995401817939 & 0.575394508856 & 0.296322943261 & $<10^{-48}$\\
$(0.35,0.6)$ & -1.266182766546 & 0.304613560249 & 0.444252224180 & $<10^{-48}$\\\bottomrule
\end{tabular}
\end{center}
The printed arctangent $J$ differs by $0.415303503064$ and
$0.343007339335$, respectively.

\paragraph{Why the original check passes.}
\Lc The original \texttt{smx\_residual\_ambiguity\_interface.py} explicitly
sets \texttt{J\_cont = pi/2 + I12\_same(a1,a2+pi)}. It does not call its
own antisymmetric \texttt{J12\_same} there and does not integrate a crossed
phase. The target offset is built into this definition. Carrying the
light-square relation $J=I+\pi/2$ into another orientation is a possible
\emph{chosen continuation convention}, not a verification of the original
antisymmetric integral. Nor are the real absolute-value and sign formulas
holomorphic expressions that can be continued without specifying branches.

\Lc Direct resolved quadrature at $\alpha_1=0.2$ and
$\alpha_2=\pi/3\pm10^{-7}$ finds the raw same-direction $I$ and $J$
continuous; the $J$ upper-minus-lower difference is about
$3.50\,10^{-7}$. No constant $+\pi/2$ jump of this raw $J$ occurs there.
Different root-sum extensions past the stability cusp may carry different
constants and must not be substituted silently for this integral.
The proof also uses the false identity
$\arctan x-\arctan(1/x)=\frac\pi2\operatorname{sgn}x$; it is the
\emph{sum}, for real nonzero $x$, that has this value. At $x=2$ the
printed difference is $0.643501108793$, not $\pi/2$.
The claimed $\cos3\alpha>0$ on the whole $(0,\pi/3)$ interval is false;
it changes sign at $\pi/6$.

\Lc Even if the stipulated continuation offset $C=\pi/2$ were adopted,
crossing $-c_M\pi/2+c_{\rm loop}\Delta_{\rm same}$ with
$\Delta_{\rm cross}+\Delta_{\rm head}=C$ would give
$(c_M-c_{\rm loop})\pi/2+c_{\rm loop}\Delta_{\rm head}$,
not $(c_M-\tfrac12c_{\rm loop})\pi/2+c_{\rm loop}\Delta_{\rm head}$.
For a regular good head-on pair with $\alpha_1+\alpha_2<\pi/3$,
$\eta_{\rm ho}<0$ and the Stirling offset is $+\pi/2$; the corollary's
parenthetical $\eta_{\rm ho}>0$ has the wrong sign.
\Ld An absolute determinant phase or Maslov cancellation cannot be inferred
from the imposed offset. The claimed physical cancellation is unestablished.

\section{Small-amplitude interface: the checker and text disagree}
\paragraph{Corrected chain-rule terms (\texttt{interfaces}).}
\Lc Use the shifts in \eqref{eq:correctwall} and put
$\alpha_1=\lambda r$, $\alpha_2=\lambda$, $0<r<1$,
$\Delta P(\alpha,+)=c_{\ln}\alpha\log\alpha+c_{\rm lin}\alpha+
O(\alpha^3\log\alpha)$. With $\mathcal C=c_{\ln}\log\lambda+c_{\rm lin}$,
\begin{align}
 H_{\rm ov}&\longrightarrow
 -2c_{\ln}\log r\log\frac{1-r}{1+r},\\
 D_{\rm ov}&=-2(r^{-1}-r)\mathcal C
             -2c_{\ln}\frac{\log r}{r}+o(1),\\
 (H+D)_{\rm ho}&=-2(r^{-1}+r)\mathcal C
             -2c_{\ln}\frac{\log r}{r}+o(1).
 \label{eq:small}
\end{align}
Here $H$ uses the physical Hirota shifts and $D$ the Darboux additions.
Our action derivatives and finite-$\lambda$ tests reproduce these formulas
at $\lambda=10^{-12}$. They agree with the \emph{original script}, not its
write-up. In particular the overtaking tail is generally logarithmically
divergent, not a finite scale-invariant constant.

\Lc The independently evaluated mode-charge formula gives
\begin{equation}
 c_{\ln}=\frac1\pi,\qquad
 c_{\rm lin}=\frac{-\log2-1+\tfrac14\log3}{\pi},
\end{equation}
not the text's $c_{\ln}=-2/\pi$. Also
$A_{\rm ho}(\lambda r,\lambda)=1+3r\lambda^2+O(\lambda^4)$,
not the printed $1-\tfrac32\lambda^2(1-r^2)$.
For the raw same-direction continuum integral,
\begin{equation}
 I_{\rm KdV}(r)=-\frac4\pi\chi_2(r)
 -\frac2\pi\log r\log\frac{1-r}{1+r},\qquad
 J_{\rm KdV}(r)=\frac\pi2+I_{\rm KdV}(r).
\end{equation}
An independent elementary quadrature of the kernel verifies these expressions;
at fixed $r\ne1$, $J$ is generally nonzero, contrary to the printed
$O(\lambda^2)$ assertion. With the actual mode coefficient, the logarithmic
product cancels in $I_{\rm KdV}-H_{\rm ov}$, not in the sum.
\Ld These remain conditional frame conversions and continuum integrals;
none replaces a projected determinant or the missing two-body frame map.

\section{Dependency, circularity, and execution audit}
\paragraph{Phase 2i dependency check.}
\Lc The package explicitly excludes Phase 2i, and its four scripts import
no Phase 2i implementation. No Wiener--Hopf coefficient, endpoint product,
or all-loop master identification from that phase is used. Thus its refuted
claims are not an upstream dependency here. The present quadrature is
consistent with the orientation bookkeeping confirmed in
\texttt{audit\_phase2i.tex}; we do not use the refuted universal kernel
criterion, printed Darboux defect coefficient, or proposed physical phase
selection. This also leaves intact that audit's warnings about physical
normalisation and the common quantum frame.

\paragraph{What the original scripts establish and omit.}
\Lc All four named scripts return exit code zero and end with
\texttt{ALL CHECKS PASSED}. This was checked both directly and by the new
checker's subprocess tests. Their shared dictionary is an implementation
dependency, not by itself circular. The following limitations matter:
\begin{itemize}
\item The $J$ crossing offset is inserted in \texttt{J\_cont}; no independent
crossed $J$ or cusp-to-Stirling phase transport is calculated.
\item The mirror check differentiates the factorised ansatz and is a valid
conditional algebra test; it cannot show that constant-$G$ products fail
unitarity. The numerical dilog fit is not its advertised global proof.
\item The pole checker does not test the false sign boundary or positivity
threshold, the printed $\hbar=0.05$ examples, or the claimed $|D|=\pi/3$
resonance. Its actual numerical pole examples use $\hbar=0.1$.
\item The ring checker assumes the seam twist in \texttt{Phi\_ring}, and
then tests its consequences. This is useful consistency testing, but does
not independently derive the twist or quantise winding. Its chamber loop
uses $N=4$ only and its sum loop $N=3$ only, despite broader docstring wording.
\item The residual-interface script uses corrected Darboux signs and
corrected small-amplitude expressions; passing them does not validate the
contradictory formulas in the PDF. It does not prove the residual logarithm
exists globally, or that the interface is unique.
\end{itemize}
\Ld No inference about authors' intentions or editing history follows from
these mismatches. Successful scripts cannot substitute for checking that
they implement the printed claims.

\paragraph{Failures recorded.}
\begin{enumerate}
\item \Lc Target-claim discrepancies explicitly asserted by the independent
checker: reversed Darboux signs; incorrect dilog residues and blanket
exclusion; false good-breather sign and positivity boundaries; incorrect pole
examples; false overtaking resonance and smoothness; omitted Bethe self
term; printed continuum $I,J$ mismatch; imposed $J$ crossing constant;
false arctangent identity and constant conversion; incorrect small-amplitude
tails, head-on expansion, and $J$ limit. Qualified physical claims are
identified separately rather than asserted false by numerical tests.
\item \Lc The first independent cusp assertion expected an upper-minus-lower
raw same-direction $J$ jump of $-\pi/2$. It failed: resolved direct quadrature
gave $3.50034825911617\,10^{-7}$ at separation $2\,10^{-7}$. Inspection of
the elementary phases and their endpoint product showed the raw same-direction
integral is continuous here. The final test asserts continuity at the same
finite-step tolerance; no tolerance was enlarged to hide this failure.
This correction strengthens the distinction between raw quadrature and a
chosen continuation of a closed-form kernel.
\item \Lc A later independent run failed a symbolic assertion because only
the first half of an identically zero expression was expanded before testing
structural equality with zero. Expanding the entire difference fixed the
checker; the identity, numerical tolerances, and physical conclusion were
unchanged. The quadratures preceding this assertion had all passed.
\end{enumerate}

\paragraph{Reproducibility and deliverables.}
\Lc Run \texttt{python3 phase2/scripts/audit\_smx.py}. Its final line is
\texttt{ALL CHECKS PASSED}; refuted statements are encoded as assertions
that their discrepancies exist. The original scripts' successful execution
is separately labelled as a smoke test. The report is compiled by two
\texttt{pdflatex} runs in \texttt{phase2/}, with rendered pages inspected.
The deliverables are \texttt{audit\_smx.tex}, \texttt{audit\_smx.pdf}, and
\texttt{scripts/audit\_smx.py}. No existing status line or research text was
changed, as required by this task.

\Ld The package is useful as a catalogue of conditional candidates and
classical consistency checks. Its corrected algebra does not yet select
an exact S-matrix, fix the two-body quantum frame, or establish the claimed
quantum ring and absolute one-loop phase interpretation.
\end{document}
